\section{Appendix D: Original-record lower bounds and sharp comparison}

The calibrated families in \cref{obj:def:continuous-calibrated-priors} and their joint support certificate in \cref{obj:lem:calibrated-support} provide the alternatives for the expected-length lower bounds. The comparisons concern sample laws on the original record space \leanref{sym:\Omega}{\ensuremath{\Omega}}, including the observed covariates. We first fix the normalization of the distance used to compare these laws.

\begin{definitionv}{synth_6}[Squared Hellinger distance for original-record sample laws]
For probability laws \(Q,Q'\) on the original-record sample space \(\Omega^n\), define their unhalved squared Hellinger distance by
\[
\mathcal H^2(Q,Q')
=
\int_{\Omega^n}
\left(
\sqrt{\frac{dQ}{dR}}
-
\sqrt{\frac{dQ'}{dR}}
\right)^2\,dR,
\]
where \(R\) is any common dominating measure. The value is independent of the choice of \(R\) and lies in \([0,2]\). For the continuous calibrated mixtures, take \(R\) to be the original-record sample reference measure. This normalization assigns coefficient one to the displayed integral.
\end{definitionv}

The normalization in \cref{obj:synth_6} fixes the numerical constants in the mixture comparisons. Both comparisons use the radius and envelope constant from the same joint certificate in \cref{obj:lem:calibrated-support}. The following bounds apply when the average number of observations per partition cell satisfies the stipulated occupancy threshold.

\begin{lemmav}{lem:original-record-mixtures}[Original-record mixture bounds]
Let \(\varepsilon\) and \(M\) be the radius and envelope constant selected from the fixed joint calibrated-support certificate, with \(0<\varepsilon\le 1/100\) and \(M\ge 1\). Define the positive constants
\[
\kappa=\frac{1}{96\exp(1)},
\qquad
C_{\mathrm{mix}}=96\exp(2)\,8^2M^4.
\]
For sample laws \(Q,Q'\), use the unhalved squared Hellinger distance
\[
\mathcal H^2(Q,Q')
=
\int
\left(
\sqrt{\frac{dQ}{dR}}-
\sqrt{\frac{dQ'}{dR}}
\right)^2\,dR,
\]
where \(R\) is the original-record sample reference measure.

For every choice satisfying
\begin{itemize}
\item \textbf{(Sample and partition sizes.)} \(n,k\) are integers with \(n\ge 1\) and \(k\ge 1\).
\item \textbf{(Occupancy.)} \(n/k\le\kappa\).
\item \textbf{(Signed amplitudes.)} \(\eta,\zeta,\delta\in\mathbb R\) satisfy
\[
|\eta|\le\varepsilon,\qquad
|\zeta|\le\varepsilon,\qquad
|\delta|\le\varepsilon,
\]
\end{itemize}
let \(Q_b\), for \(b\in\{0,1\}\), and \(Q_t\) be the uniform finite mixtures of the calibrated mixed and fair original-record sample laws, respectively:
\[
Q_b=\mathbb E_\sigma\!\left[P_{b,\sigma}^{\otimes n}\right],
\qquad
Q_t=\mathbb E_\sigma\!\left[P_{t,\sigma}^{\otimes n}\right],
\]
where the expectation averages uniformly over the finite sign vectors. Let \(P_{*,t,\delta}\) be the deterministic fair comparator obtained from the fair construction at partition size \(k\), center \(t\), and amplitude \(\delta\), using the constant numerical sign vector with every entry equal to \( -1 \). Then
\[
\mathcal H^2(Q_0,Q_1)
\le
C_{\mathrm{mix}}\frac{n^2}{k}\eta^2\zeta^2,
\]
and, for every \(t\in[0,1/4]\),
\[
\mathcal H^2\!\left(Q_t,P_{*,t,\delta}^{\otimes n}\right)
\le
C_{\mathrm{mix}}\frac{n^2}{k}\delta^4.
\]
\end{lemmav}

The two bounds in \cref{obj:lem:original-record-mixtures} quantify the statistical proximity of the same-model alternatives certified by \cref{obj:lem:calibrated-support}. The mixed comparison depends on the squared product of the propensity and outcome amplitudes; the fair comparison depends on the fourth power of its outcome amplitude and is uniform over the stated centers. Both retain the factor \(n^2/k\) and use a common constant determined by the joint density envelope. Exact singleton matching and the signed derivative bounds in \cref{obj:lem:calibrated-support} supply the calibration properties underlying these comparisons.

Expected length is evaluated on the radius slice in \cref{obj:def:radius-model}, while coverage is required throughout the ambient model by \cref{obj:def:honest-procedures}. Consequently, alternatives in the ambient model constrain interval length at laws in an evaluation slice. The parametric comparison records this implication uniformly over all radii, including zero.

\begin{lemmav}{lem:parametric-null-floor}[Parametric lower bound on length]
Let \(\alpha,\beta,r\in\mathbb R\) and \(n\in\mathbb N\) satisfy:
\begin{itemize}
\item \textbf{(Exponent domain.)} \((\alpha,\beta)\in\mathcal D\), the public exponent domain given by
\[
\mathcal D=\{(\alpha,\beta)\in\mathbb R^2:0<\beta<1/4,\ \beta<\alpha\le1\}.
\]
\item \textbf{(Sample size.)} \(n\ge1\).
\item \textbf{(Radius.)} \(0\le r\le1/2\).
\end{itemize}
Then the minimax expected interval length \(J_n(\alpha,\beta;r)\) defined in \cref{obj:def:length-objective} satisfies
\[
J_n(\alpha,\beta;r)\ge \frac{3}{16}n^{-1/2}.
\]
\end{lemmav}

At zero radius, \cref{obj:lem:parametric-null-floor} gives a positive uncertainty floor for expected length under the null slice. The coverage requirement still ranges over the ambient model, so uncertainty about neighboring effects enters the criterion even when length is evaluated at zero effect. The statement also covers procedures using the independent randomization seed \leanref{sym:U}{\ensuremath{U}}, since these procedures belong to the class defining the minimax objective in \cref{obj:def:length-objective}.

Combining this parametric floor with the mixed and fair comparisons gives the lower side of the sharp expected-length profile. The upper side is attained by the literal interval output \leanref{sym:I^{\mathrm{up}}}{\ensuremath{I^{\mathrm{up}}_{n,\mathsf E,\mathsf N}(\alpha,\beta)}} defined in \cref{obj:def:upper-handle}, with its coverage and length guarantees from \cref{obj:thm:finite-honest-upper}. The resulting comparison holds throughout the public exponent domain \leanref{sym:\mathcal D}{\ensuremath{\mathcal D}} and for every admissible engine and naming policy.

\begin{theoremv}{thm:sharp-honest-frontier}[Sharp honest interval frontier]
There exist positive finite absolute constants \(c,C\) with the following guarantees. Write
\[
\mathcal D=\{(\alpha,\beta)\in\mathbb R^2:0<\beta<1/4,\ \beta<\alpha\le1\},
\]
and define the numerator and denominator exponents and the expected-length profile by
\[
\mathsf a(\alpha,\beta)
=\min\left\{\frac12,\frac{2(\alpha+\beta)}{2\alpha+2\beta+1}\right\},
\qquad
\mathsf b(\beta)=\frac{4\beta}{4\beta+1},
\qquad
\rho_n(\alpha,\beta;r)
=n^{-\mathsf a(\alpha,\beta)}+r n^{-\mathsf b(\beta)}.
\]

For every admissible rational interval-arithmetic engine \(\mathsf E\) and every admissible Borel policy \(\mathsf N\) satisfying the quantitative contracts of \cref{obj:def:arithmetic-engine,obj:def:dyadic-name-convention} and selecting valid names for the public exponents and observed covariates, consider parameters satisfying:
\begin{itemize}
\item \textbf{(Exponent domain.)} \((\alpha,\beta)\in\mathcal D\).
\item \textbf{(Sample size.)} \(n\in\mathbb N\) and \(n\ge1\).
\item \textbf{(Radius.)} \(0\le r\le1/2\).
\end{itemize}
Let \(I^\star_{n,\mathsf E,\mathsf N}(\alpha,\beta)
=I^{\mathrm{up}}_{n,\mathsf E,\mathsf N}(\alpha,\beta)\) denote the interval returned by the fixed-budget upper construction. With \(\mathcal A_n(\alpha,\beta)\) as in \cref{obj:def:honest-procedures}, define the minimax expected length by
\[
J_n(\alpha,\beta;r)
=
\inf_{I\in\mathcal A_n(\alpha,\beta)}
\sup_{P\in\mathcal M_r(\alpha,\beta)}
\mathbb E_{P^{\otimes n}\otimes\lambda}
\bigl|I(\mathcal O,U)\bigr|,
\]
where \(\mathcal M_r(\alpha,\beta)\) is the prescribed effect-radius slice, \(\mathcal O\) is the sample of original records, and \(U\) is an independent seed with uniform law \(\lambda\) on the unit interval. Then
\[
c\rho_n(\alpha,\beta;r)
\le J_n(\alpha,\beta;r)
\le
\sup_{P\in\mathcal M_r(\alpha,\beta)}
\mathbb E_{P^{\otimes n}}
\bigl|I^\star_{n,\mathsf E,\mathsf N}(\alpha,\beta)\bigr|
\le C\rho_n(\alpha,\beta;r).
\]
The same constants apply to every admissible engine and naming policy.

For every such engine and policy, every \((\alpha,\beta)\in\mathcal D\), and every \(n\ge1\), the actual name-dependent output is a deterministic, radius-independent Borel interval with global \(90\%\) coverage over \(\mathcal M(\alpha,\beta)\), the prescribed ambient model. It is a nonempty closed connected interval contained in \([-1/2,1/2]\), with rational endpoints, returned after the two prescribed finite precision budgets and finite rational arithmetic. Its inputs consist of the fixed engine, the public exponents and their valid public names, and the \(n\) original records and their valid covariate names. Fixing the explicit admissible engine \(\mathsf E_0\) and dyadic-floor naming policy \(\mathsf N_0\) gives the concrete attaining interval
\[
I_n^\star
=I^{\mathrm{up}}_{n,\mathsf E_0,\mathsf N_0},
\qquad
c_\star=c,\qquad C_\star=C,\qquad n_0=1,
\]
where \(c_\star,C_\star\) are the comparison constants and \(n_0\) is the sample-size threshold.

Moreover, for every \((\alpha,\beta)\in\mathcal D\), every \(0\le r\le1/2\), every \(n\ge1\), and every \(I\in\mathcal A_n(\alpha,\beta)\), including procedures using independent randomization,
\[
c\rho_n(\alpha,\beta;r)
\le
\sup_{P\in\mathcal M_r(\alpha,\beta)}
\mathbb E_{P^{\otimes n}\otimes\lambda}
\bigl|I(\mathcal O,U)\bigr|.
\]

For every \((\alpha,\beta)\in\mathcal D\), the exponent difference is positive and satisfies
\[
\mathsf a(\alpha,\beta)-\mathsf b(\beta)>0,
\qquad
\mathsf a(\alpha,\beta)-\mathsf b(\beta)
=
\begin{cases}
\displaystyle
\frac{2(\alpha-\beta)}
{(2\alpha+2\beta+1)(4\beta+1)},
&\alpha+\beta\le1/2,\\[6pt]
\displaystyle
\frac{1-4\beta}{2(4\beta+1)},
&\alpha+\beta\ge1/2.
\end{cases}
\]
At \(\alpha+\beta=1/2\), the two displayed expressions agree.

Finally, for every nonempty compact set \(K\subseteq\mathcal D\), the same constants and threshold satisfy the compact-uniform conditions
\[
0<\inf_{(\alpha,\beta)\in K}c,
\qquad
\sup_{(\alpha,\beta)\in K}C<\infty,
\qquad
\sup_{(\alpha,\beta)\in K}1=1.
\]
\end{theoremv}

The two terms in \cref{obj:thm:sharp-honest-frontier} distinguish uncertainty in the numerator coordinate from effect-scaled uncertainty in the denominator coordinate. The mixed alternatives supply the radius-independent lower contribution, together with the parametric floor when the numerator exponent reaches \(1/2\); the fair alternatives supply the contribution proportional to the evaluation radius. Thus the null slice retains the numerator rate, whereas every fixed positive radius eventually makes the slower denominator rate dominant. The positive exponent difference stated in the theorem establishes this ordering throughout the prescribed domain. A single globally honest interval attains the comparison simultaneously across radii.

The following synthesis collects the statistical comparison, the finite-budget realization from \cref{obj:lem:fixed-budget-realization}, and the concrete engine supplied by \cref{obj:lem:concrete-arithmetic-engine}. It identifies the attaining procedure with the actual output for each admissible engine and naming policy.

For the concrete sequence below, the inputs are the sample size, public exponent representations, and represented original records, together with the independently fixed arithmetic engine and naming policy. The same radius-independent sequence is evaluated over each radius subclass.

\begin{theoremv}{thm:full-honest-frontier-solution}[Full honest frontier solution]
There exist absolute constants \(c,C\in\mathbb R\), with \(0<c\) and \(0<C\), for which the following assertions hold. Write
\[
\mathcal D=\{(\alpha,\beta)\in\mathbb R^2:0<\beta<1/4,\ \beta<\alpha\le1\},
\qquad
\mathsf a(\alpha,\beta)=
\min\left\{\frac12,\frac{2(\alpha+\beta)}{2\alpha+2\beta+1}\right\},
\qquad
\mathsf b(\beta)=\frac{4\beta}{4\beta+1},
\]
and define the expected-length profile by
\[
\rho_n(\alpha,\beta;r)
=d_n(\alpha,\beta;r)
=n^{-\mathsf a(\alpha,\beta)}
+r n^{-\mathsf b(\beta)}.
\]

For every admissible rational interval-arithmetic engine \(\mathsf E\), every admissible Borel naming policy \(\mathsf N\), both satisfying the quantitative contracts of \cref{obj:def:arithmetic-engine,obj:def:dyadic-name-convention}, every \((\alpha,\beta)\in\mathcal D\), and every integer \(n\ge1\), take
\[
I^\star_{n,\mathsf E,\mathsf N}(\alpha,\beta)
=I^{\mathrm{up}}_{n,\mathsf E,\mathsf N}(\alpha,\beta),
\]
the literal output of the fixed-budget represented upper construction. This procedure belongs to the globally \(90\%\)-honest Borel class \(\mathcal A_n(\alpha,\beta)\) of \cref{obj:def:honest-procedures}. It receives neither the radius nor unknown nuisance functions and terminates on every valid named input within the prescribed rank and endpoint budgets of \cref{obj:def:resolutions,obj:def:upper-handle,obj:lem:fixed-budget-realization}. For every sample and randomization seed its output is a closed interval \([l,u]\) with
\[
l,u\in\mathbb Q,
\qquad -\frac12\le l\le u\le\frac12.
\]
For every \(r\in[0,1/2]\), the same constants satisfy
\[
c\rho_n(\alpha,\beta;r)
\le J_n(\alpha,\beta;r)
\le
\sup_{P\in\mathcal M_r(\alpha,\beta)}
E_P\!\left[
\left|I^\star_{n,\mathsf E,\mathsf N}(\alpha,\beta)(\mathcal O,U)\right|
\right]
\le C\rho_n(\alpha,\beta;r).
\]
Here \(J_n(\alpha,\beta;r)\) is the minimax expected length of \cref{obj:def:length-objective}, and the expectation uses \(n\) original records drawn independently from \(P\), together with the independent randomization seed \(U\).

In particular, the explicit concrete engine \(\mathsf E_0\) and the dyadic-floor naming policy \(\mathsf N_0\) give the sequence
\[
I_n^\star(\alpha,\beta)
=I^{\mathrm{up}}_{n,\mathsf E_0,\mathsf N_0}(\alpha,\beta).
\]
For every \((\alpha,\beta)\in\mathcal D\) and \(n\ge1\), this sequence is globally honest, has the rational closed ordered output described above at every input, and satisfies the same displayed comparison for every \(r\in[0,1/2]\).

Moreover, for every \((\alpha,\beta)\in\mathcal D\), every \(r\in[0,1/2]\), every integer \(n\ge1\), and every \(I\in\mathcal A_n(\alpha,\beta)\), including procedures using independent randomization,
\[
c\rho_n(\alpha,\beta;r)
\le
\sup_{P\in\mathcal M_r(\alpha,\beta)}
E_P\!\left[|I(\mathcal O,U)|\right].
\]

The exponent difference is positive throughout \(\mathcal D\) and obeys
\[
\mathsf a(\alpha,\beta)-\mathsf b(\beta)=
\begin{cases}
\displaystyle
\frac{2(\alpha-\beta)}
{(2\alpha+2\beta+1)(4\beta+1)},
&\alpha+\beta\le1/2,\\[6pt]
\displaystyle
\frac{1-4\beta}{2(4\beta+1)},
&\alpha+\beta\ge1/2.
\end{cases}
\]
At \(\alpha+\beta=1/2\), the two displayed expressions agree.

Finally, select the lower and upper comparison constants and sample-size threshold as
\[
c_\star(\alpha,\beta)=c,\qquad
C_\star(\alpha,\beta)=C,\qquad
n_0(\alpha,\beta)=1.
\]
For every nonempty compact \(\mathcal K\subset\mathcal D\), these choices satisfy
\[
0<
\inf_{(\alpha,\beta)\in\mathcal K}c_\star(\alpha,\beta),
\qquad
\sup_{(\alpha,\beta)\in\mathcal K}C_\star(\alpha,\beta)<\infty,
\qquad
\sup_{(\alpha,\beta)\in\mathcal K}n_0(\alpha,\beta)=1.
\]
\end{theoremv}

The statistical and computational guarantees in \cref{obj:thm:full-honest-frontier-solution} concern the same interval: its rational endpoint construction preserves global coverage and the sharp expected-length comparison at every sample size. The original-record lower bound covers the entire honest class, including randomized procedures, while each admissible implementation attains the comparison with the same absolute constants. Selecting these constants and the threshold \(n_0=1\) also gives uniform comparisons on every nonempty compact subset \leanref{sym:\mathcal K}{\ensuremath{\mathcal K}} of the public exponent domain.
